\subsection{Combined controls and broader use of the available observations}
\label{sec:available_data}
% Generated from completed saved comparisons.
\newcommand{\JointCoreOneDelta}{23.69}
\newcommand{\JointCoreThreeDelta}{6.70}
\newcommand{\GasFortyDelta}{18.87}

We extend the conventional checks through combined covariance/core exclusions, a different gas architecture, individual ESPRESSO exposures and additional archived absorbers. These extensions follow inspection of the original target and are exploratory; the new-target screen uses declared quality criteria. Rules and intermediate endpoints are retained in the analysis records.

The combined controls apply the preceding frozen 2382/2600 core masks and tapered correlation kernel simultaneously. Covariance uses retained original pixel indices; both hypotheses re-optimize the same 45-component gas model and continua. Excluding model flux below 0.1 leaves 2692 pixels and gives $\Delta\chi^2=\JointCoreOneDelta$ for five added shifts. The wider exclusion, below 0.3 with three-pixel padding on either side, leaves 2473 pixels and gives $\Delta\chi^2=\JointCoreThreeDelta$ (Table~\ref{tab:available_data}). Improvements from different pixel selections cannot be combined into one likelihood-ratio statistic.

Masking also removes sensitivity. At the fixed diagonal-noise $H_0$ solution of Sec.~\ref{sec:espresso_null}, we project the shift derivatives against gas and continuum derivatives under the GLS weights. For the illustrative direction in which 2382 and 2600 each shift by $100\,\mathrm{m\,s^{-1}}$ while the other three relative shifts remain zero, the mild and wide masks retain approximately 70.5\% and 16.0\% of the full-pixel local information. This calculation ignores nuisance-parameter bounds and characterizes that single direction, not the marginal information for five free shifts. Thus the weaker wide-mask improvement cannot alone establish that saturation caused the original offsets.

\begin{figure}[htbp]
\centering
\includegraphics[width=.96\textwidth,height=.34\textheight,keepaspectratio]{joint_controls_information.pdf}
\caption{Combined controls and information removed by masking. Left: nonlinear improvements under diagonal and empirical GLS errors, each adding five shifts. Right: retained local information for one specified coherent strong-line shift, after nuisance projection at the fixed preceding baseline. This is conditional sensitivity, not evidence for a detected signal or all information in the data.}
\label{fig:joint_controls_information}
\end{figure}

We also replace the published 45-component architecture with a fixed 40-component alternative. Adjacent published component centres separated by less than $1.025\,\mathrm{km\,s^{-1}}$, half the smallest adopted instrumental FWHM, are joined into connected groups. Their initial total column density, column-weighted velocity and Doppler-core second moment are preserved; this does not preserve full Voigt profiles or assert that sub-resolution components are unmeasurable. Both hypotheses fit all 2931 pixels with the same empirical covariance, using 132 and 137 parameters. The original bound widths are centred on the merged initialization. The five-shift improvement is $\GasFortyDelta$, and the 2382 and 2600 shifts become approximately $-89$ and $-80\,\mathrm{m\,s^{-1}}$, changing by about $22\,\mathrm{m\,s^{-1}}$ from the 45-component GLS result. Two tested alternative-model starts terminate at objectives differing by $12.95$, so the selected local fit retains substantial initialization uncertainty. This tests sensitivity to a specified architecture and its bounds, not a nested statistical selection of the true number of clouds.

\begin{table}[htbp]
\caption{Combined masking/covariance controls and a gas-structure sensitivity test. Here $k_0$ counts the parameters of the conventional hypothesis. Each within-row comparison adds five relative shifts, using the same pixels, covariance and nuisance freedom. The empirical covariance retains the supplied variance amplitude. Results are bounded local fits; solver stopping criteria do not certify global minima. Cross-row objective differences are not independent discovery statistics.}
\label{tab:available_data}
\centering
\small
\begin{tabular}{lrrrr}
\hline
GLS comparison & Pixels & $k_0$ & $\chi^2_0$ & $\Delta\chi^2$ \\
\hline
Full pixels, 45 components & 2931 & 147 & 2231.33 & 29.24 \\
Joint $F<0.1$ mask + GLS & 2692 & 147 & 1931.90 & 23.69 \\
Joint padded $F<0.3$ mask + GLS & 2473 & 147 & 1684.36 & 6.70 \\
Full pixels, 40 components & 2931 & 132 & 2255.16 & 18.87 \\
\hline
\end{tabular}

\end{table}

For an exposure-level check, we retrieve all 17 extracted ESPRESSO S2D spectra and fit their native vacuum-barycentric pixel bins. The published exposure-specific rejections and atmospheric masks are transferred conservatively; no second barycentric correction is applied. This diagnostic fixes the gas template from the preceding combined-spectrum baseline while fitting line velocities and local continua and zero levels. The pooled shifts relative to 2374 are $-41.9\pm16.4$ and $-45.5\pm15.0\,\mathrm{m\,s^{-1}}$ for 2382 and 2600. Applying the same estimator to the combined spectrum gives $-42.8$ and $-39.6\,\mathrm{m\,s^{-1}}$. Agreement here checks two representations of the same observations; a full gas re-fit uses a different estimator. The shared data-derived template precludes an independent confirmation, and the quoted errors condition on that template and diagonal extraction errors.

The same 2600 transition appears in adjacent echelle orders. Its lower-index minus upper-index velocity is $-60.2\pm16.7\,\mathrm{m\,s^{-1}}$, or $-57.7\pm16.6\,\mathrm{m\,s^{-1}}$ when each order has a fitted Gaussian width; some free-width solutions reach their bounds. An order-dependent difference in the same physical transition directs attention to calibration, extraction, instrumental response and template approximations. It does not show that these effects explain the entire original offset, and its post-inspection conditional uncertainty is not a calibrated discovery test. Exposure-year comparisons likewise concern observational consistency, not a measurable change in the absorber's cosmic age.

Further UVES fits use expected-fluctuation errors and widen the Gaussian-width scale bounds from $[0.65,1.25]$ to $[0.5,1.4]$. The two strong-line shifts change by only $4$--$5\,\mathrm{m\,s^{-1}}$, while the 2374 width scale reaches $0.638$, below its former boundary. However, all four 1800-evaluation continuations and a 400-evaluation null cross-start fail the independently checked first-order stationarity criteria. Their finite objective differences therefore receive no calibrated significance. This symmetric-width control leaves covariance, calibration, non-Gaussian instrumental response and gas-model uncertainty unresolved; the independent photons still share gas architecture and atomic inputs with ESPRESSO.

The archive extension screens 155 identified damped Ly$\alpha$ absorbers along 132 sightlines within the 467-spectrum SQUAD archive \cite{Murphy2019SQUAD}. A coverage, forest, atmospheric-band and catalogue signal-to-noise screen selects 36 absorbers along 32 distinct sightlines under the initial atmospheric policy; broader water-band exclusions leave seven metadata candidates. All 32 spectra were retrieved and all 324 nine-transition windows analyzed. Six absorbers at $1.629\leq z\leq2.659$ pass the subsequent automatic coverage, noise and multiplet-detection criteria; three retain absorption-at-window-edge warnings. These are candidates for further modelling, not six calibrated precision measurements. The DLA catalogue is a targeted subset, not an exhaustive search for every Fe\,II absorber in 467 spectra, and systems on the same sightline share calibration.

We perform bounded Voigt-profile pilots for two compact candidate systems, selecting displayed component counts by exploratory $H_0$ AICc among 1, 2, 3, 4 and 6 components. In J232128$-$105122 at $z=1.629$, three shared gas components fitted to five transitions give conventional $\chi^2=80.59$ for nominal $\nu=101$; four additional shifts improve it by $6.84$. In J004131$-$493611 at $z=2.248$, six components and three transitions give $\chi^2=113.19$ for $\nu=108$, with an improvement of $5.66$ for two shifts. Both use the supplied statistical-error array with a diagonal covariance approximation; re-fitting with the archive's expected-fluctuation array gives improvements of $6.16$ and $5.50$, respectively. Component count and initialization matter: placing one shared component initially near a weak feature resolves an initially poor fit without increasing the number of components. These are initial model comparisons, not calibrated constant-variation measurements.

The 32 additional reduction logs contain no explicit nonzero exposure shift/slope entries; this does not prove that every calibration correction is absent. Wavelength distortions and gas structure require further assessment. Shared templates, atomic references and calibration assumptions also prevent treating all diagnostics as independent evidence. These extensions supply neither a calibrated $\Delta\alpha/\alpha$ measurement nor a homogeneous sample with a specified physical response for testing a cosmic-time law.
